\section{Discussion}
The main outcome of the thesis is not that one architecture wins every
comparison, but that the ranking depends on the level at which the models are
judged. The matched next-token evaluation gives the clearest result: under the
chosen tokenizer, context length, batch size, and update budget, xLSTM is the
strongest language model. The generated-MIDI evaluation is less one-dimensional.
The descriptor-level metrics favor different models for different musical
properties, and the FMD ordering does not exactly follow the held-out
next-token ordering. This difference is important because it shows that
validation loss is useful for checkpoint selection, but it is not a complete
proxy for musical generation quality.

The xLSTM result suggests that the recurrent architecture is well matched to the
available data and token representation. With only a 2,000-file subset, the
model may benefit from a stronger sequential inductive bias and from a parameter
scale that is large enough to model the REMI+ stream without requiring the same
amount of data as a more general Transformer baseline. The medium xLSTM
follow-up also shows the boundary of this conclusion: scaling capacity and
training time together did not automatically improve the validation objective.
For this dataset size, checkpoint selection and regularization appear at least
as important as raw model size.

The Transformer baselines are still informative even though they do not lead on
held-out next-token metrics. They provide a controlled reference point for how
decoder-only self-attention behaves under the same symbolic-music setup. The
LLaMA-style model is the stronger Transformer on token prediction and pitch
distribution agreement, while the GPT-2-style model gives the lowest FMD in the
final generation comparison. This means the Transformer results should not be
treated as failed baselines. They show that different architectural choices can
preserve different statistical properties of the generated MIDI, even when their
language-model loss is worse.

The strongest practical tradeoff is between predictive accuracy and generation
cost. xLSTM gives the best held-out token performance and strong rhythmic
density agreement, but its autoregressive generation is much slower in the final
run. That cost is plausible for a recurrent model with sLSTM blocks because each
new token depends on updating recurrent state step by step. The Transformer
models are weaker in the matched token metrics but substantially faster during
generation. For an offline evaluation study, the slower xLSTM path is
acceptable; for an interactive tool, the speed difference matters more.

The evaluation setup therefore supports a cautious model-choice rule rather
than a single universal winner. If the priority is next-token modelling of the
held-out REMI+ stream, the selected xLSTM checkpoint is the best choice. If the
priority is fast interactive sampling, the Transformer models are more practical
despite weaker token metrics. If the priority is a particular musical statistic,
the choice depends on the statistic: rhythmic sparsity, pitch distribution, and
FMD do not all select the same model.

This also affects how the Gradio applications should be interpreted. They are
not another quantitative benchmark; they are a way to inspect model behavior at
the level where aggregate metrics are least complete. MIDI playback, iterative
continuation, and generation diagnostics make it possible to notice problems
such as sparse note output, unnatural repetition, or poor continuation behavior
that are only partly captured by the automatic metrics. In that sense, the apps
are a useful bridge between the stored experiment artifacts and later listening
or user studies.

The results should still be read within the limits of the experimental design.
The comparison is controlled on the most important practical variables, but it
is not an exhaustive architecture search. The Transformer configurations are
matched baselines, not separately optimized systems, and the xLSTM conclusion is
tied to the chosen tokenizer, subset size, schedule, and sampling protocol. A
larger dataset, a broader Transformer sweep, or a generation protocol with hard
duration limits could change the balance between models. The current evidence is
therefore strongest as a controlled comparison of three implemented pipelines
under one shared setup, and weaker as a general claim about xLSTM or
Transformer architectures for symbolic music in all settings.
